% =========================================================
% Rationals Chapter Roadmap
% =========================================================

\section*{Chapter Roadmap --- The Rational Numbers}
\label{sec:rationals-chapter-roadmap}

This chapter develops the rational numbers in stages.  At each stage, the new structure is introduced and then the theorems proving that the structure works are placed immediately after it.

\begin{tcolorbox}[colback=gray!8, colframe=gray!55!black, arc=2pt,
  left=6pt, right=6pt, top=4pt, bottom=4pt,
  title={\small\textbf{Toolkit --- Rational Number Development}},
  fonttitle=\small\bfseries]
\[
\text{equivalence classes}
\longrightarrow
\text{arithmetic}
\longrightarrow
\text{field}
\longrightarrow
\text{order}
\longrightarrow
\text{ordered field}
\longrightarrow
\text{approximation}
\longrightarrow
\text{gaps}
\longrightarrow
\mathbb{R}.
\]
\end{tcolorbox}

\begin{enumerate}
  \item Construct rational numbers as equivalence classes of integer pairs.
  \item Define arithmetic using representatives and prove well-definedness.
  \item Prove the algebraic laws of rational arithmetic.
  \item Package those laws by proving that $\mathbb{Q}$ is a field.
  \item Define positivity and order on $\mathbb{Q}$.
  \item Prove that $\mathbb{Q}$ is an ordered field.
  \item Develop linear approximation: density, dyadics, small denominators, mediants, Farey sequences, and Stern--Brocot.
  \item Develop nonlinear approximation: square bracketing and rational square approximations.
  \item Convert repeated approximation into rational sequences and Cauchy behavior.
  \item Exhibit rational gaps, especially the square-root gap at $2$.
  \item Transition to the real numbers as the completion of rational arithmetic and order.
\end{enumerate}
